\section{A common cohomological operator}\label{sec:translation}

We now use the logarithmic families to turn an elementary modification of Higgs bundles into monodromy of Betti spaces. All cohomology identifications in this section are the restriction maps of Corollary~\ref{log:restriction}. This specifies a single operator before we estimate its effect on either filtration.

\subsection{Modification of two eigenlines}

Choose $v\in\sqrt{-1}\mathbb R^n$ with sum zero, distinct entries, and no vanishing proper nonempty partial sum. Such vectors exist because the excluded conditions form finitely many proper real hyperplanes. Set
\[
\ell=(1,-1,0,\ldots,0).
\]
At a logarithmic operator with distinct residue eigenvalues, a positive
elementary modification allows a simple pole in a chosen residue eigenline;
a negative one takes the kernel of the projection to that eigenline at $x$.
With the convention that a positive shift replaces a local generator $s$
by $z^{-1}s$, the $\lambda$-Leibniz rule changes its residue eigenvalue by
$-\lambda$. A negative shift changes it by $+\lambda$.
We perform the shifts specified by $\ell$ successively: first the positive
modification on the first labeled eigenline, then, after reselecting the
labeled eigenspaces of the new residue, the negative modification on the
second labeled eigenline. The ordered eigenspaces of the final residue
define the final flag.

\begin{lemma}\label{trans:translation}
The two modifications define a continuous family of isomorphisms
\begin{equation}\label{trans:path}
F_a:Y_{a,v}\longrightarrow Y_{a,v-a\ell},\qquad 0\le a\le1.
\end{equation}
At $a=0$ they define an algebraic automorphism of $Y_{0,v}$. More generally, the same modification defines an automorphism over the punctured Higgs line $(\lambda,u)=(0,zv)$, $z\ne0$, preserving the logarithmic Hitchin map. At $a=1$, the Riemann--Hilbert identifications make $F_1$ the identity on the underlying Betti space.
\end{lemma}

\begin{proof}
For each single modification, choose a local frame adapted to its current
residue eigenspaces. The off-diagonal coefficients involving the chosen
eigenline vanish at the pole, so the singly modified lattice remains
preserved by the logarithmic operator. The preceding calculation gives its
new residue eigenvalues. Their imaginary parts remain distinct throughout
\eqref{trans:path}, including after the first step. We can therefore
reselect the residue eigenspaces before performing the second modification.
The two-step composite preserves the determinant and its operator because
the lattice shifts have sum zero. Its inverse performs the opposite
single modifications in reverse order, again reselecting eigenspaces
between steps. These elementary transformations and their inverses vary
algebraically where the successive residue eigenvalues are distinct,
giving the asserted continuous family.

We check stability, which is not supplied by lattice modification alone. On the Higgs line with $z\ne0$, a proper invariant saturated subbundle would have a trace Higgs differential with a single possible simple pole. Its residue would be $z\sum_{i\in I}v_i$ for a proper nonempty subset $I$. This is nonzero, contrary to the residue theorem. The modified Higgs bundles satisfy the same argument. Distinct residue eigenvalues split the local Higgs operator into formal eigenline summands; shifting these summands and shifting them back proves invertibility.

For $a>0$, divide the operator by $a$. A proper invariant monodromy subspace would have determinant monodromy one around the puncture. The imaginary part of the sum of its exponents is $a^{-1}\sum_{i\in I}\operatorname{Im}v_i$, which is nonzero. Thus no such subspace exists. Differences of exponents have nonzero imaginary part, so no nonzero integral resonance occurs. The inverse modifications are valid and remain stable by the same tests.

The characteristic polynomial on the Higgs line is unchanged away from the pole and hence unchanged as a meromorphic polynomial on $C$. Finally, modifications of a connection change only its logarithmic lattice at the puncture. At $a=1$ the exponent changes are integers, so the meromorphic flat bundle and the ordered monodromy eigenlines are unchanged. This is precisely the endpoint identification asserted in the lemma.
\end{proof}

\subsection{Translation as eigenvalue monodromy}

Write $T=F_0$ and transport its pullback to the total family:
\begin{equation}\label{trans:global-T}
\mathsf T=\rho_{0,v}^{-1}T^*\rho_{0,v}
\quad\text{on }H^*(\mathcal Y,\Q).
\end{equation}
To compare this operator with the endpoint $F_1$, we need constancy along
the two families in \eqref{trans:path}. Their fibers are nonproper; the
restriction isomorphisms, rather than properness, supply the following
criterion.

\begin{lemma}
\label{log:compact-cycles}
Let $I$ be a connected real interval and let $p:M\to I$ and
$q:N\to I$ be smooth submersions, with finite-dimensional rational
cohomology on their fibers.  Suppose there are spaces $A,B$ and
continuous maps $a:M\to A$, $b:N\to B$ whose restrictions induce
isomorphisms
\[
 \alpha_t:H^*(A,\mathbb Q)\xrightarrow{\sim}H^*(M_t,\mathbb Q),
 \qquad
 \beta_t:H^*(B,\mathbb Q)\xrightarrow{\sim}H^*(N_t,\mathbb Q)
\]
for every $t\in I$.  If $F:M\to N$ is continuous over $I$, then
$\alpha_t^{-1}F_t^*\beta_t:H^*(B,\mathbb Q)\to H^*(A,\mathbb Q)$
is independent of $t$.
\end{lemma}

\begin{proof}
Fix a degree and a parameter $t_0$.  Choose finitely many compact
singular cycles forming a rational homology basis of $M_{t_0}$.
A lift of the interval vector field near the compact union of their
carriers has a flow for a common short time.  It transports these
cycles to nearby fibers; the resulting cycles sweep homologies in
$M$.  Their images in $A$ are therefore homologous, so their
pairings with the restrictions of a fixed basis of $H^*(A)$ are
constant.  The restriction isomorphisms show that the transported
cycles remain a basis of the fiber homology.

Applying the continuous map $bF$ to these swept cycles gives
homologies in $B$.  Hence their pairings with every fixed class
of $H^*(B)$ are also constant.  These pairings are the matrix
coefficients of $\alpha_t^{-1}F_t^*\beta_t$, relative to the
fixed bases.  The operator is thus locally constant, including
one-sided neighborhoods of interval endpoints, and therefore
constant on $I$.
\end{proof}

Apply the lemma with $I=[0,1]$ and with $M,N$ the pullbacks of
$\mathcal Y\to\Lambda$ along
\[
 a\longmapsto(a,v),\qquad a\longmapsto(a,v-a\ell),
\]
respectively. These are smooth submersions by
Proposition~\ref{log:families}, and \eqref{trans:path} gives the continuous
map $F:M\to N$. Take both ambient spaces in the lemma to be
$\mathcal Y$, with the natural maps from the two pullbacks. Their fiber
restrictions induce the isomorphisms of
Corollary~\ref{log:restriction}. The constancy conclusion therefore gives
\begin{equation}\label{trans:endpoint}
\mathsf T=\rho_{1,v}^{-1}F_1^*\rho_{1,v-\ell}.
\end{equation}
In particular, this identification does not require the nonproper smooth families to be globally locally trivial.

Recall the eigenvalue map
\[
t(u)=(\zeta e^{-2\pi\sqrt{-1}u_i})_{i=1}^n.
\]
The path $u=v-a\ell$, $0\le a\le1$, maps to a loop in $V$. Every forbidden proper subproduct has modulus different from one because the corresponding partial sum of $v$ has nonzero imaginary part. The same condition excludes resonance along the entire path. By Proposition~\ref{log:rh}, this is a path in the actual moving-puncture Betti family. Restrictions of classes from $\mathcal Y$ are flat sections along this path: the Betti direct images satisfy ordinary base change and are locally constant. Since $F_1$ is the endpoint identity, \eqref{trans:endpoint} identifies $\mathsf T$ with the pullback convention for its eigenvalue monodromy.

\begin{lemma}\label{trans:logarithm}
There is an integer $M>0$ such that
\begin{equation}\label{trans:N}
N=\frac1M\log(\mathsf T^M)
\end{equation}
is a well-defined nilpotent derivation of $H^*(\mathcal Y,\Q)$. Let
$N_{\lambda,u}=\rho_{\lambda,u}N\rho_{\lambda,u}^{-1}$.
On $H^*(Y_{1,0},\Q)$, equipped with the diagram weights of Section~\ref{sec:diagram}, one has
\[
N_{1,0}W_j\subset W_{j-2}.
\]
\end{lemma}

\begin{proof}
The full-flag Betti family has integral lisse direct images, ordinary base change, and Hodge--Tate fibers by Section~\ref{sec:betti}. The moving-puncture transition maps are algebraic and locally constant by Proposition~\ref{log:rh}. The finite-monodromy statement of Section~\ref{sec:diagram} therefore applies to this very family. Choose $M$ that kills the monodromy action on every weight-graded piece of its total cohomology. There are finitely many such pieces. Then $\mathsf T^M-1$ is nilpotent, and \eqref{trans:N} is a finite logarithmic sum. The weights are even, so $N$ lowers them by at least two at $(1,v)$.

The path $(1,av)$, $0\le a\le1$, stays in the nonresonant inverse image of $V$. At $a=0$, the scalar point belongs to $V$ because $\zeta$ has order $n$. For $a>0$, the imaginary-part argument applies again. Parallel transport preserves the diagram weights and coincides with the restriction identifications. This gives the stated estimate at $(1,0)$.

Lemma~\ref{diag:finite} makes the logarithm a derivation on the moving-puncture Betti cohomology. The restriction identifications are ring isomorphisms, so this property transports to $H^*(\mathcal Y,\Q)$.
\end{proof}

\subsection{Flag push-pull and the weight bound}

For the scalar correspondence at $\lambda$, put
\[
\mathsf A_\lambda=\pi_{\lambda,*}\iota_\lambda^*,
\qquad
\mathsf B_\lambda=\iota_{\lambda,*}\pi_\lambda^*.
\]
Its unflagged moving-pole space is $S\times X_\lambda$, with projection $\pr_\lambda$. We define
\begin{equation}\label{trans:fprime}
f'=\pr_{0,*}\,\mathsf A_0 N_{0,0}^2\mathsf B_0\,\pr_0^*
:H^m(X,\Q)\longrightarrow H^{m-2}(X,\Q).
\end{equation}
The flag embedding and projection have opposite degree shifts, $2b$ and
$-2b$. Integration over the compact curve $S$ supplies the remaining
degree $-2$. Table~\ref{spec:shift-table} records the individual derived morphisms
and their degree and weight shifts. Section~\ref{sec:specialization}
derives the perverse bound from their composite over the Hitchin base.
We first prove the weight bound.

\begin{proposition}\label{trans:weights}
Under nonabelian Hodge theory, the operator \eqref{trans:fprime} satisfies
\[
f'W_jH^m(X^B,\Q)\subset W_{j-4}H^{m-2}(X^B,\Q).
\]
\end{proposition}

\begin{proof}
The scalar correspondence and its Gysin maps commute with the restriction identifications of Corollary~\ref{log:restriction}. Thus \eqref{trans:fprime} transports to the same expression at $\lambda=1$. By Proposition~\ref{log:nah}, these identifications transport the ordinary Betti weight filtration on the full cohomology of $X^B$. The possible difference from a chosen nonabelian-Hodge identification is one algebraic $\Gamma$-action, which preserves every weight subspace, including those in variant cohomology.

At $\lambda=1$, the scalar Betti space is the flat algebraic product $S\times X^B$ of Lemma~\ref{log:scalar-flat-product}. Its diagram mixed Hodge structure assigns weight zero to the cochains of $S$ and the ordinary Deligne weights to $H^*(X^B)$. Hence both $\pr_1^*$ and $\pr_{1,*}$ preserve weights. The flag maps are maps of the same flat transition diagrams. Their ordinary topological Gysin maps are the diagram Gysin maps, with weight shifts $2b$ for $\mathsf B_1$ and $-2b$ for $\mathsf A_1$. These shifts cancel, while $N_{1,0}^2$ lowers weight by four by Lemma~\ref{trans:logarithm}.
\end{proof}

We have constructed the operator and proved its weight bound. The remaining tasks are to realize it over the Hitchin base, which gives the perverse bound, and to compute its double commutator with $e$.
